\begin{abstract}
This paper proposes an auditable ``static ontology--protocol readout'' framework to formally characterize how three-dimensional spatial representation, temporal sequence, and material stability emerge as products of finite-resolution readout from a single high-dimensional information structure. The ontological layer posits a static six-dimensional integer lattice point structure as a microscopic phase--index backbone, and avoids directly equating it with empirical space. The protocol layer centers on cut-and-project and scan--readout: observation is equivalent to selecting a three-dimensional observable subspace and its orthogonal complement equipped with a compact acceptance window, and performing discrete scans along the complementary phase parameters to generate quasi-periodic point sets on a three-dimensional screen and update readout sequences with each scan iteration. To suppress periodic locking and enhance distinguishability, this paper formulates the ``golden sector'' as a choice of minimal resonance irrational orientation/scan slope, whose mathematical characterization is compatible with minimal partial quotients of continued fractions, Sturmian sequences, and Zeckendorf/Ostrowski encoding. Furthermore, under finite-window readout, the paper introduces stable-type sets $X_m^{\mathrm Z}$ and folding maps $\Fold_m$ to objectify indistinguishable degrees of freedom as preimage fibers $\mathcal{F}_m$; the temporal arrow is defined at the protocol layer as a contraction mechanism of posterior constraint sets/cylinder measures, and ``future'' strictly corresponds to the uncertainty magnitude of conditional distributions (such as conditional entropy). Finally, the paper presents a closed-loop model of the adaptive observer: the readout operator is controlled by internal parameters and updated by data-driven observations, forming a self-referential protocol dynamical system. The closing section lists testable propositions, covering quasi-periodic spectral lines, discrete degenerate structures of resolution folding, protocol-reconstruction explanations of delayed-choice-like experiments, and light-cone boundary readout interfaces under causal structure constraints.
\end{abstract}

\begin{sloppypar}
\textbf{Keywords:} Cut-and-project; quasi-periodic point sets; scan readout; finite resolution; Zeckendorf/Ostrowski encoding; stable sectors; self-referential observation; light-cone boundary readout.
\end{sloppypar}

\tableofcontents
\newpage
